% Quelle: 5. Übung CT Lösungen.pdf
% Art: Übung mit Lösungen
% Themen: Lineare Codes, Hamming-Code, Hamming-Schranke, Singleton-Schranke, 1-dimensionale Unterräume
\section{Übung 5: Wissen aus den Lösungen}
% Die Quelle ist ein Übungsblatt mit Lösungen. Eigenständige Definitionen/Sätze enthält sie nicht.
% Die folgenden Aussagen werden in den Lösungen (Seiten 3--5) wörtlich verwendet und hier nur gesammelt;
% der vollständige Kontext steht in uebungen/05_uebung.tex.

\subsection{In den Lösungen verwendete Aussagen}

\begin{bemerkung}[Kontrollmatrix von $H_q(m)$, aus Lösung 5\_5]
Die Spalten einer Kontrollmatrix $H$ des $q$-nären Hamming Codes $H_q(m)$ bestehen aus Repräsentanten der 1-dimensionalen Unterräume des $\mathbf{F}_q^m$.
% KORRIGIERT: Voraussetzung h \neq 0 fehlte (für h = 0 ist die Menge {0} kein 1-dim. Unterraum und hat nicht q Elemente)
1-dim.\ Unterraum des $\mathbf{F}_q^m$ sind die Geraden (mit $\mathbf{h} \in \mathbf{F}_q^m$, $\mathbf{h} \neq \mathbf{0}$)
\[
\{\, y \cdot \mathbf{h} \mid y \in \mathbf{F}_q \,\}, \qquad |\{\, y \cdot \mathbf{h} \mid y \in \mathbf{F}_q \,\}| = |\mathbf{F}_q| = q .
\]
\end{bemerkung}

\begin{bemerkung}[Parameter von $H_q(m)$, aus Lösung 5\_5b]
\begin{align*}
\operatorname{Rang} G + \operatorname{Rang} H &= n = (q^m - 1)/(q-1)\\
\operatorname{Rang} H = m, \quad \operatorname{Rang} G &= k = n - m\\
\Rightarrow \quad |C| &= q^k
\end{align*}
\end{bemerkung}

\begin{bemerkung}[Hamming Schranke, aus Lösung 5\_6]
% KORRIGIERT: stand nur "2e+1 = d" ohne Voraussetzungen; das gilt nur für ungerades d, allgemein e = \lfloor (d-1)/2 \rfloor
Für einen Code $C \subseteq \mathbf{F}_q^n$ mit Minimalabstand $d$ und $e = \lfloor (d-1)/2 \rfloor$ gilt
\[
|C| \cdot \sum_{i=0}^{e} \binom{n}{i} (q-1)^i \le |\mathbf{F}_q^n| = q^n, \qquad 2e+1 = d \text{ (für ungerades } d\text{)} .
\]
\end{bemerkung}

\begin{bemerkung}[Singleton Schranke, aus Lösung 5\_7]
% KORRIGIERT: Voraussetzung (linearer [n, k, d]_q-Code) fehlte
Für einen linearen $[n, k, d]_q$-Code gilt
\[
k + d \le n + 1 .
\]
\end{bemerkung}
